\section{Method}
\label{sec:method}

\subsection{Problem Setting and Notation}

We consider binary classification. Let $x_i \in \mathbb{R}^d$ be the feature
vector of individual $i$, $y_i \in \{0,1\}$ the true label, and
$s_i \in \mathcal{S}$ the sensitive attribute, where $\mathcal{S}$ is a finite
set of demographic groups. A model $f_\theta$ produces a logit
$z_i = f_\theta(x_i)$, a probability $p_i = \sigma(z_i)$ with
$\sigma(\cdot)$ the logistic function, and a hard decision
$\hat{y}_i = \mathbf{1}[p_i \geq 0.5]$.

\subsection{Why a Surrogate Is Necessary}

The evaluation criteria of Section~\ref{sec:setup} are defined on the hard
decisions $\hat{y}_i$. The map $p_i \mapsto \mathbf{1}[p_i \geq 0.5]$ has zero
derivative almost everywhere and is undefined at the threshold, so any
objective written in terms of $\hat{y}$ has no usable gradient and cannot be
minimized by backpropagation. We therefore construct \emph{differentiable
surrogates} that operate on the probabilities $p_i$ directly, and verify
empirically (Section~\ref{sec:surrogate}) that each surrogate tracks the hard
metric it stands in for.

\subsection{Group Term: Soft Demographic Parity}

The demographic parity difference is the largest gap in selection rate between
any two groups. Replacing the selection rate with the mean predicted
probability gives a differentiable analogue:
\begin{equation}
\softdp(p, s) \;=\; \max_{g \in \mathcal{S}} \bar{p}_g \;-\; \min_{g \in \mathcal{S}} \bar{p}_g ,
\qquad
\bar{p}_g \;=\; \frac{1}{|\mathcal{I}_g|}\sum_{i \in \mathcal{I}_g} p_i ,
\label{eq:softdp}
\end{equation}
where $\mathcal{I}_g = \{ i : s_i = g \}$. Equation~\eqref{eq:softdp} is
differentiable in $\theta$ through $p$, and converges to the hard demographic
parity difference as the probabilities saturate towards $\{0,1\}$.

\subsection{Individual Term: Soft Generalized Entropy}

Following Speicher et al.~\cite{speicher2018unified} we define the
\emph{benefit} received by individual $i$ as
\begin{equation}
b_i \;=\; 1 + p_i - y_i .
\label{eq:benefit}
\end{equation}
A correctly served individual has $b_i \approx 1$; a false positive receives
more than deserved ($b_i \to 2$) and a false negative less ($b_i \to 0$).
Individual unfairness is then the \emph{inequality} of the benefit vector,
measured by the generalized entropy index of order $\gamma$:
\begin{equation}
\softge(p,y) \;=\; \frac{1}{n\,\gamma(\gamma-1)} \sum_{i=1}^{n}
\left[ \left( \frac{b_i}{\bar{b}} \right)^{\gamma} - 1 \right],
\label{eq:softge}
\end{equation}
where $\bar{b} = n^{-1}\sum_i b_i$ is the mean benefit.
We use $\gamma = 2$ throughout, for which \eqref{eq:softge} equals half the
squared coefficient of variation of $b$. The index is zero exactly when every
individual receives identical benefit, and grows as the burden of the model's
errors concentrates on fewer people. Because $b_i$ depends on $p_i$ smoothly,
\eqref{eq:softge} is differentiable.

\subsection{The Composite Loss}

The two fairness terms are combined with the accuracy term into a single
objective with two interpolation weights $\alpha, \beta \in [0,1]$:
\begin{equation}
\boxed{\;
\mathcal{L}(\theta)
= \alpha\,\mathcal{L}_{\mathrm{BCE}}
+ (1-\alpha)\Big[\beta\,\softdp + (1-\beta)\,\softge\Big] \;}
\label{eq:composite}
\end{equation}
where $\mathcal{L}_{\mathrm{BCE}}$ is binary cross-entropy computed from
logits for numerical stability. The two parameters have distinct roles:

\begin{itemize}
\item $\alpha$ trades \emph{accuracy against fairness}. At $\alpha=1$ the loss
      reduces to plain cross-entropy, giving the unconstrained baseline. As
      $\alpha \to 0$ accuracy is ignored entirely, which yields degenerate
      constant predictors (Section~\ref{sec:selection}).
\item $\beta$ trades \emph{group against individual fairness}, and is the
      variable of interest in this study. At $\beta=1$ the fairness budget is
      spent entirely on demographic parity; at $\beta=0$ entirely on benefit
      equality; intermediate values blend them.
\end{itemize}

Sweeping $\beta$ at fixed $\alpha$ therefore holds the total fairness pressure
constant while redistributing it between the two notions---which is exactly
the comparison needed to isolate a trade-off between them.

\subsection{Models}

We instantiate \eqref{eq:composite} with two model families, chosen to span a
linear and a nonlinear hypothesis class while remaining standard in the
fairness literature:
\begin{itemize}
\item \textbf{Logistic regression}: a single affine layer returning a logit.
\item \textbf{MLP (64--64)}: two hidden layers of 64 units with ReLU
      activations, returning a logit.
\end{itemize}
Both return logits rather than probabilities; the sigmoid is applied once,
inside the loss.

\subsection{Training and Selection Procedure}
\label{sec:selection}

Algorithm~\ref{alg:sweep} states the full procedure. Training uses
Adam~\cite{kingma2015adam} with learning rate $0.01$ and weight decay
$10^{-4}$, for at most $300$ epochs, with early stopping on the validation
value of the \emph{same} composite loss being optimized (patience $20$) and
restoration of the best weights from a deep copy. Because \eqref{eq:softdp}
requires all groups to be present in order to form a group gap, each epoch is
a single full-batch gradient step rather than a pass of mini-batches.

\begin{algorithm}[t]
\caption{$\alpha\times\beta$ sweep with validation-only selection}
\label{alg:sweep}
\begin{algorithmic}[1]
\Require dataset $D$, seed set $S$, grids $A$ (for $\alpha$), $B$ (for $\beta$)
\For{each seed $k \in S$}
  \State $(D_{\text{tr}},D_{\text{va}},D_{\text{te}}) \gets \textproc{Split}(D, 60/20/20, k)$
  \State fit scaler and encoder on $D_{\text{tr}}$; apply to all splits
  \For{each architecture $a$ and each $(\alpha,\beta) \in A \times B$}
     \State seed weights with $k$; initialise $f_\theta$
     \State train $f_\theta$ on $D_{\text{tr}}$ minimising \eqref{eq:composite}
     \State early-stop on $\mathcal{L}$ evaluated on $D_{\text{va}}$
     \State record test metrics, and validation accuracy,
            \dpdtwo{} and positive rate
  \EndFor
  \State \textbf{// selection uses validation quantities only}
  \State $F \gets \{(\alpha,\beta) : \alpha < 1,\;
         0.05 \le \widehat{r}_{\text{va}} \le 0.95 \}$
  \State $(\alpha^\ast,\beta^\ast) \gets
         \arg\min_{F} \sqrt{(1-\mathrm{acc}_{\text{va}})^2
         + \mathrm{DPD}_{2,\text{va}}^{2}}$
  \State report test metrics at $(\alpha^\ast,\beta^\ast)$
\EndFor
\State aggregate over $S$; paired Wilcoxon test with Holm correction
\end{algorithmic}
\end{algorithm}

Two details in Algorithm~\ref{alg:sweep} matter for the validity of the
results. First, \emph{test metrics are never used for selection}: the
$(\alpha,\beta)$ pair is chosen by distance to the utopia point
$(\text{accuracy}=1, \dpdtwo=0)$ computed on validation data alone. Second,
the constraint $0.05 \le \widehat{r}_{\text{va}} \le 0.95$ on the validation
positive rate $\widehat{r}_{\text{va}}$ excludes near-constant predictors: a
model that assigns every individual to one class attains
$\dpdtwo = 0$ trivially while being useless. This guard is not cosmetic---at
$\alpha=0$ we observe validation positive rates above $0.99$.

\subsection{Pipeline Overview}

Fig.~\ref{fig:pipeline} summarizes the complete experimental pipeline from raw
data to reported results.

\begin{figure}[t]
\centering
\begin{tikzpicture}[
  node distance=3.2mm,
  every node/.style={font=\scriptsize},
  box/.style={rectangle, rounded corners=1.5pt, draw=black!70, thick,
              text width=0.80\columnwidth, align=center, inner sep=2.4pt,
              minimum height=6.2mm, fill=black!3},
  hi/.style={box, fill=blue!8, draw=blue!55},
  sm/.style={rectangle, rounded corners=1.5pt, draw=black!70,
             text width=0.375\columnwidth, align=center, inner sep=2.4pt,
             minimum height=8.5mm, fill=black!3},
  ar/.style={-{Latex[length=1.6mm]}, thick, draw=black!70}
]
\node[box] (data) {\textbf{Data}: COMPAS ($n{=}6{,}172$), German ($n{=}1{,}000$),
                   Adult ($n{=}45{,}222$)};
\node[hi, below=of data] (filt) {\textbf{Soundness filters}: ProPublica row
                   filters; \emph{drop label-leaking} \texttt{duration}};
\node[box, below=of filt] (split) {\textbf{Per-seed split} $60/20/20$,
                   seeds $0$--$9$};
\node[box, below=of split] (prep) {\textbf{Preprocessing} (fit on train only):
                   StandardScaler $+$ one-hot};
\node[sm] (sweep) at ([xshift=-0.205\columnwidth, yshift=-11mm]prep.south)
      {$\boldsymbol{\alpha\times\beta}$ \textbf{sweep}\\
       $7\times7$, 2 archs\\ $2{,}940$ models};
\node[sm] (base)  at ([xshift= 0.205\columnwidth, yshift=-11mm]prep.south)
      {\textbf{Baselines}\\ Reweighing, ExpGrad,\\ ThreshOpt, RF, XGB};
\node[hi] (sel) at ([yshift=-24mm]prep.south)
                  {\textbf{Selection on validation only}
                   $+$ degeneracy guard};
\node[box, below=of sel] (eval) {\textbf{Test evaluation}: Acc, \dpdtwo{},
                   EOD, \getwo{}};
\node[box, below=of eval] (agg) {\textbf{Aggregate over 10 seeds};
                   Wilcoxon $+$ Holm};
\node[hi, below=of agg] (res) {\textbf{Trade-off map}:
                   accuracy--fairness frontiers, $\beta$ cross-effect curves,
                   master and ablation tables};

\draw[ar] (data) -- (filt);
\draw[ar] (filt) -- (split);
\draw[ar] (split) -- (prep);
\draw[ar] (prep.south) -- ++(0,-3.2mm) -| (sweep.north);
\draw[ar] (prep.south) -- ++(0,-3.2mm) -| (base.north);
\draw[ar] (sweep.south) |- ([yshift=-3.2mm]sweep.south) -| (sel.north);
\draw[ar] (base.south)  |- ([yshift=-3.2mm]base.south)  -| (sel.north);
\draw[ar] (sel) -- (eval);
\draw[ar] (eval) -- (agg);
\draw[ar] (agg) -- (res);
\end{tikzpicture}
\caption{Experimental pipeline. Shaded boxes mark the three steps that most
affect the validity of the conclusions: removal of the label-leaking COMPAS
feature, selection performed strictly on validation data with a
degenerate-predictor guard, and reporting of the resulting trade-off map.}
\label{fig:pipeline}
\end{figure}
